\newpage

\ifdefined\useChapters
\chapter{A Porta Vermelha e o Vazio}
\else
\chapter{}
\fi

Já faz uma semana que eu despertei, e é um porre: todo mundo me trata como se eu pudesse quebrar a qualquer momento. Eu não vou, gente, não vou. É estranho dormir e acordar adolescente. Às vezes soa como brincadeira de mau gosto, mas ninguém aparece pra mostrar a câmera e encerrar a ilusão.

Algumas coisas mudaram desde que abri os olhos. Meus pais não estavam mais vivos e, junto com eles, perdi a casa que eu chamava de lar. Os conhecidos ainda existem, mas a convivência perdeu cor. Sempre achei piegas quando diziam que sentimento é planta que precisa de água; hoje entendo. As pessoas são gentis, me fazem companhia — e mesmo assim algo dentro de mim ficou fora do lugar.

Não consigo parar de pensar no que vivi durante o coma. Dizem que foi só invenção da minha mente. Eu não acredito. Por que eu escolheria ficar preso com um cara que fingia ser minha esposa? Por que sentiria medo de algo sem nome, e ao mesmo tempo um acolhimento que doía de tão real? Talvez eu queira que tenha sido de verdade porque ali, de algum modo, eu pertencia.

Não dá mais pra viver sem cor. Preciso juntar as lembranças, decifrar o que der pra decifrar. Não lembro de tudo, mas alguns detalhes grudaram: a casa amarela, antiga, dois andares; o jardim descuidado; uma árvore meio seca; crianças gritando ao longe; gente falando português com o meu sotaque. Peças demais, e eu sem manual.

Os dias passaram e a esperança foi ficando fina. A escola piorou: antes eu sofria por ser invisível, agora me olham como o “cara mais velho” que “inventou histórias”. O novo jogo é apontar e sussurrar. Voltar pra casa é o melhor do dia — lá eu só preciso encarar a Lúcia.

Um fim de tarde, no ônibus, uma pedra bateu no vidro e o motorista reduziu. Olhei pela janela e vi. A casa. A bendita casa. Amarela — mas viva, jardim aparado, janelas sem tábuas. Um arrepio: a mesma, e não a mesma.
 
— Meu Deus!
— Para esse ônibus, motorista!
— Para! Eu preciso descer aqui!

O motorista freou. As pessoas me fitaram, metade raiva, metade curiosidade. Eu já estava na calçada quando ouvi o ronco do motor recomeçando.

Corri até o portão. Nada de árvore seca, nem clima de terror. Luz boa, cheiro de grama molhada. Bati três vezes. Eu sempre bato três vezes quando estou nervoso.

A porta abriu. Ela — olhos de mel, mais linda do que a memória. Fiquei mudo por um instante.

— Pois não, posso te ajudar? — perguntou, desconfiada, me medindo dos pés à cabeça.

— Você mora aqui faz tempo?

— Oi? Se você veio vender alguma coisa, não viu a placa no muro? — disse, já fechando a porta.

Não era a pessoa doce da lembrança. Ou talvez a vida tenha arrancado isso dela. Segurei a porta com o pé sem pensar.

— Desculpa. Procurei essa casa por muito tempo. Eu... eu estive aqui. Em coma. Ou achei que estive.

Ela suspirou, reabriu a porta um palmo, depois mais um.

— Você sabe ganhar atenção, rapaz. Entra e fala logo.

Por dentro, a lareira, as cadeiras, a escada para os quartos — tudo como na minha cabeça, só que mais nítido, com reformas discretas. Contei o que podia sem soar maluco. O nome do cachorro, Quincas. O nome do homem: Crispim. O cheiro de fumaça, o peso das cortinas. Ela ouviu em silêncio, tomou um gole de chá e disse:

— Vem comigo.

Caminhamos até uma sala de portas vermelhas — as mesmas que, no sonho, eu não devia abrir. Lá dentro, móveis antigos e um senhor sentado numa cadeira de balanço. Cabeça baixa. Barbas e cabelos grisalhos descuidados. No colo, um cachorro dormindo.

— Meu pai, Crispim — disse ela, chorando. — Está catatônico há mais de quinze anos. Uma vez ou outra se agita e solta uma palavra, mas... é isso.

Era ele. Muito mais velho, mas ele. Meu estômago apertou.

— Por que você veio? — ela perguntou, a voz falhando. — Pra zombar da minha cara?

— Não. Não ganho nada com isso. Eu não sei o que está acontecendo. Só sei que, de alguma forma, estive aqui.

— Então por ora você está me perturbando. Quero paz. Vai embora.

— Como eu saberia dessa sala? Como eu saberia o nome do seu pai?

— Por favor. Vai.

Saí. A certeza e a dúvida brigavam dentro de mim. Eu poderia estar louco. Ou poderia estar certo demais. Voltei pra casa com mais perguntas e uma convicção: se eu acordei, talvez pudesse ajudar Crispim a acordar também.

Tic, tac, tic, tac.

O relógio de parede tem sido inimigo do sono — e talvez tenha sido ele quem me puxou de volta da primeira vez. Dias depois, quando já tinha desistido de procurar sinais, lá estava eu outra vez diante da casa amarela, agora com o jardim seco e as cortinas fechadas. Entrei sem bater. A porta vermelha me esperava como quem já me conhece.

Abri.

O mundo apagou.

\medskip

Escuro. Frio. Uma sala sem paredes, sem chão — só eu. Andar não mudava nada; o vazio se movia comigo. O som de uma xícara pousando em pires, depois uma voz:

— Sente-se. Chá?

Ele surgiu do escuro. O mesmo homem, mais novo do que na cadeira de balanço. O olhar dele prendia. Uma parte de mim ardia de raiva: e se foi ele quem me roubou cinco anos de vida? A outra parte queria respostas.

— Você deve achar que te prendi aqui — disse. — Não te culpo por ter raiva. Mas me diz: o que você faria se estivesse sozinho por tempo demais e, de repente, aparecesse a chance de não estar?

— Chamar isso de “companhia” é bonito. Mas é egoísta — respondi, duro.

Ele ficou calado, soprando o chá. Cinco minutos que pareceram pedra. Eu também calei. Às vezes o silêncio revela mais do que confissão.

— Você acha que aqui moral importa? — disse, por fim. — Já viu o que fazem lá fora? Gente que prende e tortura porque pensa diferente. Aqui, pelo menos, ninguém me encosta a mão.

A ditadura ainda vivia dentro dele. Um país preso num homem. E se ele soubesse? E se ele soubesse e escolhesse ficar?

Resolvi furar a bolha. Errei na dose.

— De onde vem a comida que não estraga? Por que ninguém bate à porta? Por que você não envelhece? — falei de uma vez.

Ele recuou um passo, apertou o roupão.

— Vou te contar — disse, cabisbaixo. — Eu descobri que podia sair do meu corpo. No começo era só isso: deixar o corpo e andar por aí. Depois aprendi a convencer a mim mesmo de outra realidade, uma onde as coisas doíam menos. Toda vez que algo incomodava, eu vinha pra cá. E é tão fácil ficar quando tudo dói muito, que um dia não consegui mais voltar.

Eu queria duvidar, e parte de mim acreditou. Enquanto ele falava, alguma coisa dentro de mim se esticava como elástico velho. Um rumor de que eu mesmo poderia fazer o mesmo, ou já estivesse fazendo.

— E a Joana? — perguntei. — Por que você usou a imagem dela pra me manter aqui?

Ele fechou os olhos um instante.

— Eu precisava de alguém. De um rosto que me lembrasse que eu existo. Pra mim, ela foi isso. Pensei que pra você também pudesse ser. Eu estava desesperado.

— Você usou a memória de quem amava pra me prender. — Minha voz tremeu.

— Talvez tenha sido cruel. Mas anos aqui dentro acabam com a gente. Quando você apareceu, achei que fosse minha chance de não estar mais sozinho.

O chão mexeu. Ou eu. As bordas da sala vibraram, como se a realidade fosse filme antigo fora de foco.

— O que você está fazendo? — Crispim arregalou os olhos. — Não tem o direito de mexer no meu mundo!

As cadeiras mudaram de lugar sozinhas; a lareira acendeu sem fogo. Quincas miou em algum canto que não existia. Eu não sabia se era eu ou ele. Ou os dois.

— Então você sabe — eu disse. — E ainda assim ficou.

Ele me olhou por um tempo que não sei medir.

— Se puder sair, saia — disse, enfim. — Se não por mim, por você. Não deixe que a minha desgraça vire a sua. Você ainda tem uma vida.

Tic, tac, tic, tac.

O som voltou primeiro. Depois a luz. Eu estava inteiro e não estava. Entendi menos do que queria; entendi o suficiente. Se eu podia entrar, eu podia sair. Se eu podia sair, talvez pudesse puxá-lo.

Quando abri os olhos, o teto era o de sempre. O relógio marcava uma hora qualquer. Eu sabia duas coisas: que minha vida não seria a mesma — e que eu ia voltar àquela porta vermelha até ela me deixar passar.

\medskip

Antes de dormir, coloquei o relógio dentro da gaveta. Cinco minutos depois, levantei e o tirei de lá. Não sabia se Crispim havia me puxado, se eu o encontrara por acaso ou se agora eu mesmo procurava a porta. Só sabia que já estava escolhendo o próximo retorno enquanto fingia ter medo dele. Deixei o relógio sobre a mesa, ao alcance da mão. Quando o tique-taque recomeçou, não me afastei.
